\section{Serious objections}

\subsection{``Scripture is sufficient; later definitions add human doctrines''}

The objection rightly fears a human authority placed above God's word and
tests whether later formulas have apostolic warrant. The Catholic answer is not
that Scripture is deficient material. Scripture is inspired and normative
within the one apostolic deposit. The dispute concerns whether Christ left that
deposit to self-interpreting individuals or to a visible apostolic Church
whose teaching office serves it. A legitimate definition identifies and
defends what was given; a purported novelty contrary to the deposit would lack
authority. Development therefore requires continuity of object, not merely an
institutional signature.

\subsection{``Later papal definitions bypass the ancient conciliar Church''}

Eastern Orthodox criticism often joins doctrinal, historical, and
ecclesiological concerns; it should not be reduced to dislike of Mary or of
certainty. Catholic teaching answers that the bishop of Rome's defining office
belongs within, not outside, the Church's apostolic faith and episcopal
communion. Consultation and reception are theologically significant even
though they are not causes that make a valid ex cathedra definition
irreformable. The remaining disagreement about primacy and the conditions of
ecclesial reception is real and cannot be dissolved by relabeling it.

\subsection{``History shows change, so irreformable truth is impossible''}

History does show changing vocabulary, questions, evidence, and conceptual
frameworks. \emph{Mysterium Ecclesiae} rejects the inference that formulas
therefore cease to signify determinate truth. A formula can be limited without
being false, and later expression can deepen without reversing its object.
Historical criticism serves faith when it distinguishes the ancient substance
from later terminology; it undermines inquiry when it assumes in advance that
continuity is impossible.

\subsection{``If teaching is not infallible, conscience is sovereign''}

Conscience is the proximate judgment one must follow, but it must be formed; it
is not a private source of revelation or immunity from ecclesial obligation.
Authentic teaching has a claim before agreement. At the same time, ``obey''
cannot erase the difference between doctrine, law, prudence, and an
administrator's factual mistake. A properly formed response seeks truth,
receives authority, states difficulty honestly, and refuses both self-will and
false certainty.

\subsection{``If obedience is required, every Roman statement is effectively
infallible''}

This maximalism protects neither Rome nor faith. It converts changeable
discipline and prudential acts into alleged revelations, then makes every
correction look like contradiction. Catholic obedience is differentiated
because ecclesial authority itself is differentiated. Precision is not
minimalism; it is justice to the authority actually exercised.

